% Set 2008
% Autors: Adelaida Clavaguera/Leopold Palomo
% Presentació de l'aplicatiu Bulmages
%
% Aquest document està alliberat sota llicència
% Creative Commons 
% http://creativecommons.org/licenses/by-sa/2.5/es/
% Licencia Reconocimiento-Compartir bajo la misma licencia 2.5 España 
% License de Creative Commons

\documentclass{beamer}
\setbeamertemplate{navigation symbols}{}
\usepackage[utf8x]{inputenc}
\usepackage[catalan]{babel}
\usepackage{eurosym}
\usepackage{url}
%\usepackage[T1]{fontenc}

\usetheme{Warsaw}

\beamersetuncovermixins{\opaqueness<1>{25}}{\opaqueness<2->{15}}

\title{BulmaGés}
\subtitle{Gestió empresarial amb programari lliure}
\author{Adelaida Clavaguera - Leopold Palomo}
\date{18 de setembre de 2008}


\institute{\begin{large}Associació Iglues\end{large} \\
 Volcànica.cat LAN Party 2008 \\ Les Preses }

% Show current subsection at subsection begining
\AtBeginSection[]{
    \begin{frame}
        \frametitle{}
        \tableofcontents[currentsubsection,hideothersubsections]
    \end{frame}
}

\begin{document}


% Title page %%%%%%%%%%%%%%%%%%%%%%%%%%%%%%%%%%%%%%%%%%%%%%%%%%%%%%%%%%%%%%%%%%%


\begin{frame}
\titlepage
\begin{overprint}
\includegraphics[scale=0.1]{iglues.eps}
\end{overprint}

\end{frame}

\begin{frame}\frametitle{Continguts de la presentació}
\tableofcontents[hideallsubsections]
\end{frame} 


\section{Presentació} 
\subsection{Qui som}
\begin{frame}\frametitle{Qui som} 
\begin{itemize}
\item Adelaida Clavaguera Mora
\item Leopold Palomo Avellaneda
\end{itemize} 
Membres de l'associació Iglues
\end{frame}

\subsection{El context actual}
\begin{frame}\frametitle{Aplicacions de gestió empresarial} 
\begin{itemize}
\item n'hi ha des de l'origen de la informàtica
\item \dots i sobretot des del PC
\end{itemize} 
Existeixen diverses alternatives al mercat, per a les PIMES
\begin{itemize}
\item Conegudes: Contaplus, Navision, SAP, Vaan \dots
\item Creades a mida
\end{itemize} 
\end{frame}

\subsection{El programari lliure}
\begin{frame}\frametitle{Què és el programari lliure} 
Què és el programari?
\begin{itemize}
\item Ordinadors només entenen llenguatge binari
\item Els humans fem servir lenguatges de programació entenibles i després
es tradueixen a binari
\item És molt difícil i quasi imposible treballar amb llenguatge binari
\end{itemize} 
\end{frame}

\begin{frame}\frametitle{Què és el programari lliure} 
1984 Sorgeix el projecte GNU per proporcionar un \textit{complete Unix-compatible software system}
\begin{itemize}
\item Es creen les bases del concepte programari lliure i la llicència GPL
\item Es defineixen les quatre llibertats: 
\begin{itemize}
 \item Utilitzar-lo com vulgui i on vulgui
 \item Redistribuir-lo a qui vulgui pels mitjans que vulgui
 \item Modificar-lo (millora o adaptació)
 \item Redistribuir les modificacions
\end{itemize} \end{itemize} 
\centering{\textbf{Programari lliure $\neq $ programari gratuït}}
\end{frame}



\section{Bulma i Iglues} 
\subsection{Bulma}
\begin{frame}\frametitle{Bulma} 
Bergantells Usuaris de GNU/Linux de Mallorca i Afegitons / Bisoños Usuarios de GNU/Linux de Mallorca y Alrededores
\begin{itemize}
\item Un grup d'usuaris de programari lliure de les Illes Balears
\item Va néixer l'any 1998/1999
\item http://bulma.net/
\end{itemize}
\end{frame}


\begin{frame}\frametitle{Bulma \dots BulmaGés} 
Sorgeix un grup de gent que vol fer quelcom per evitar la manca de programes en el sector empresarial
\begin{itemize}
 \item Trobades setmanals ``socials'' i de ``feina''
\item El programa de comptabilitat és el primer a abordar degut al marc estatal de les aplicacions de comptabilitat
\end{itemize}
\end{frame}

\subsection{Iglues}
\begin{frame}\frametitle{Iglues} 
\textbf{I}niciativa de \textbf{G}estió \textbf{L}liure \textbf{U}niversal per a \textbf{E}mpreses i \textbf{S}ocietats

\begin{exampleblock}{}
Associació sense ànim de lucre formada per desenvolupadors i usuaris 
amb l'objectiu comú de promoure i desenvolupar el programari empresarial obert i lliure
\end{exampleblock}

\begin{exampleblock}{Objectiu}
Donar marc legal de cobertura a BulmaGés
\end{exampleblock}
\end{frame}

\subsection{Els projectes}
\begin{frame}\frametitle{Projectes} 
\begin{columns}
\begin{column}{3cm}
\begin{itemize}
\item<1-> Suite BulmaGés
\item<2-> GestiONG
\item<3-> Galopin
\end{itemize}
\vspace{1cm} 
\end{column}
\begin{column}{5cm}
\begin{overprint}
\includegraphics<1>[scale=0.5]{SelectorBulmages.eps}
\includegraphics<2>[scale=0.3]{gestiong.eps}
\end{overprint}
\end{column}
\end{columns}
\end{frame}

\section{Un exemple de migració}
\subsection{El cas de l'ADENC}
\begin{frame}\frametitle{ADENC} 
\begin{block}{}
Associació per la Defensa i l'Estudi de la Natura
\end{block}
\begin{itemize}
 \item 800 socis 
    \begin{itemize}
        \item Junta directiva $\Rightarrow$ Oficina tècnica de 5 treballadors
    \end{itemize}
 \item Pressupost anual $\sim$ 250K\euro
    \begin{itemize}
        \item Declarada d'utilitat pública $\Rightarrow$ Presenta comptes anuals a l'administració
    \end{itemize}
\end{itemize}
\end{frame}

\subsection{Procés de migració}
\begin{frame}\frametitle{Migració a PL} 
\begin{itemize}
\item Migració de Macintosh a PC i portar-nos la comptabilitat: \\
  $\Rightarrow$resoldre la comptabilitat era un requisit per migrar a Linux
\item Jornades de Programari Lliure (2004)
\item Formació en comptabilitat
\item BulmaGés 0.5.1: versió operativa però no enllestida
\item Co\lgem aboració amb el projecte (intercanvi entre entitats): \\
  $\Rightarrow$ per millorar el projecte (reportar errors, proposar millores...)\\
  $\Rightarrow$ adaptació al nou pla comptable
\item Ara s'ha alliberat la nova versió que funciona:\\
  $\Rightarrow$ BulmaGés 0.11.1
\end{itemize}
\end{frame}

\section{Perquè fer servir BulmaGés}
\subsection{Aspectes previs}
\begin{frame}\frametitle{Aspectes previs} 
\begin{itemize}
\item Costos: \\
 $\Rightarrow$ instal·lar, aprendre \dots posar en marxa \\
 $\Rightarrow$ adaptar a l'empresa \\
 $\Rightarrow$ migrar les dades \\
 $\Rightarrow$ mantenir
\item Requisits: \\
 $\Rightarrow$ Sistemes Operatius: Macintosh, GNU/Linux \\
 $\Rightarrow$ Es pot utilitzar en MSWindows (no encara \dots )
\end{itemize}
\end{frame}

\subsection{PIME usuària}
\begin{frame}\frametitle{Com a PIME usuària} 
\begin{itemize}
\item Les dades comptables no estan captives
\item Independència del proveïdor (hi ha mercat)
\item Adaptabilitat
\item Dóna la possibilitat d'actualitzar el programa encara que qui l'ha fet, l'abandoni
\end{itemize}
\end{frame}

\subsection{PIME gestora comptable}
\begin{frame}\frametitle{Com a PIME gestora comptable} 
\begin{itemize}
\item Multiempresa i multiusuari
\item Es pot administrar remotament i via web
\item Reducció de costos
\item Es pot oferir al client un servei d'assesorament i adaptació del programa (informàtic)
\end{itemize}
\end{frame}

\subsection{PIME informàtica}
\begin{frame}\frametitle{Com a PIME informàtica} 
\begin{itemize}
\item Aplicació amb una tecnologia solvent
\item Es pot adaptar el programa
\item Existència d'una comunitat d'usuaris que testen i reporten errors
\item Camp de negoci no menystenible
\end{itemize}
\end{frame}

\section{Utilitzar BulmaGés}
\subsection{BulmaCont}
\begin{frame}\frametitle{Com funciona BulmaCont: COMPTABILITAT} 
\begin{itemize}
\item Multiplataforma, multiempresa, multiusuari 
\item Funcionament en xarxa (interna i externa)
\item Plans comptables: antics, PGC 2008, Pimes, Associacions sense lucre
      \dots i el que calgui
\item Nombre dígits dels comptes (des de 7 fins a 12)
\item Generació d'informes en PDF, Full de càlcul \dots
\item Corrector: localitza errors, assentaments oberts, desquadres
\item Centres de cost
\item Assentaments inte\lgem igents (predeterminats)
\end{itemize}
\end{frame}

\begin{frame}\frametitle{Com funciona BulmaCont: COMPTABILITAT} 
\includegraphics<1>[scale=0.5]{bulmacont.eps}
\end{frame}


\subsection{BulmaFact}
\begin{frame}\frametitle{Com funciona BulmaFact: FACTURACIÓ} 
\includegraphics<1>[scale=0.5]{bulmafact.eps}
\end{frame}


\subsection{BulmaTPV}
\begin{frame}\frametitle{Com funciona BulmaTPV: TERMINAL PUNT DE VENDA} 
\includegraphics<1>[scale=0.5]{bulmatpv.eps}
\end{frame}


\subsection{Informació addicional}
\begin{frame}\frametitle{Informació addicional} 
\begin{block}{Manuals d'usuari}
\url{http://www.iglues.org/wiki1/doku.php}
\end{block}

\begin{block}{Documentació tècnica}
\url{http://www.iglues.org/wiki/doku.php?id=documentacion:start}
\end{block}
\end{frame}

\section{Projectes de futur}
\begin{frame}\frametitle{Projectes de futur} 
\begin{itemize}
 \item Difusió usuaris/co\lgem aboradors \\
  $\Rightarrow$  Revista Todo linux
 \item Ampliació de les funcionalitats \\
  $\Rightarrow$ Beca per desenvolupar gestió de socis/clients d'una entitat
 \item Empreses que donin serveis a Catalunya  \\
  $\Rightarrow$ a Mallorca són dues i van saturades
\end{itemize}
\end{frame}

\begin{frame}\frametitle{BulmaGés} 
Gestió empresarial amb programari lliure
\vspace{1cm}
\begin{block}{Web}
\url{http://bulmages.org} \\
\url{http://iglues.org}
\end{block}

\begin{block}{Llista de correu}
bulmages-main@lists.berlios.de
\end{block}
\end{frame}

\end{document}

